
\titledquestion{TA's Letter Board}

During a CS101 recitation, the TAs prepare a small game using magnetic
tiles. All tiles are placed on a narrow rail, so the only allowed
operation is to swap two \textbf{adjacent} tiles.

Each tile has a distinct label from $1$ to $n$.

The TAs want to transform the current arrangement
\[
A=\langle a_1,a_2,\ldots,a_n\rangle
\]
into a target arrangement
\[
B=\langle b_1,b_2,\ldots,b_n\rangle,
\]
where both $A$ and $B$ are permutations of $\{1,2,\ldots,n\}$.

For example,
\[
\langle1,2,3\rangle
\rightarrow
\langle2,1,3\rangle
\rightarrow
\langle2,3,1\rangle
\]
uses two adjacent swaps.

The TAs first try the obvious simulation, but soon realize that $n$ may be
very large. Your task is to help them compute the minimum number of adjacent
swaps efficiently.

{

\setlength{\leftskip}{-2em}
\textbf{I.} We first convert the tile arrangement into an integer sequence.

For each tile $x$, let $\operatorname{pos}_B(x)$ denote its position in $B$.
Then construct the sequence
\[
P=\langle p_1,p_2,\ldots,p_n\rangle,
\qquad
p_i=\operatorname{pos}_B(a_i).
\]

}

\begin{parts}

%---------------------------------------------------------
\part[2]

Let
\[
A=\langle1,2,3,4,5\rangle,
\qquad
B=\langle4,1,3,5,2\rangle.
\]

Construct the corresponding integer sequence
\[
P=\langle p_1,p_2,\ldots,p_n\rangle.
\]

\begin{solution}

$p_1 = \operatorname{pos}_B(1) = 2$
$p_2 = \operatorname{pos}_B(2) = 5$
$p_3 = \operatorname{pos}_B(3) = 3$
$p_4 = \operatorname{pos}_B(4) = 1$
$p_5 = \operatorname{pos}_B(5) = 4$

$$
P = \langle 2, 5, 3, 1, 4 \rangle
$$
\end{solution}

%---------------------------------------------------------
\part[3]

For the example from part (a),
\[
P=\langle2,5,3,1,4\rangle,
\]
count the number of inversions in $P$ and list all the inversions.


\begin{solution}

$$(2, 1)$$
$$(5, 3)$$
$$(5, 1)$$
$$(5, 4)$$
$$(3, 1)$$

Totally, there are 5 inversions in $P$.

\end{solution}


%---------------------------------------------------------
\part[10]

Design a divide-and-conquer algorithm to compute the minimum number of
adjacent swaps needed to transform $A$ into $B$.

Your algorithm should use the sequence $P$ defined above and should run
faster than $\Theta(n^2)$.

Explain why your algorithm is correct.

\textbf{Hint:} Recall merge sort and inversions.

\begin{solution}

Transforming $A$ into $B$ using adjacent swaps is equivalent to sorting the sequence $P$.

Swapping the adjacent elements in $A$ corresponds to swapping the adjacent elements in $P$.

Every adjacent swap reduces the number of inversions in $P$ by exactly 1. Therefore, the minimum number of adjacent swaps needed to transform $A$ into $B$ is equal to the number of inversions in $P$.

Design the algorithm as follows:

\begin{enumerate}
    \item[]   0. Initialize: Set $l=1$, $r=n$.
    \item[]   1. Divide: $\text{mid} = \lfloor (l + r) / 2 \rfloor$, divide $[l ,r]$ into $[l, \text{mid}]$ and $[\text{mid}+1, r]$.
    \item[]   2. Conquer: Recursively count inversions in $[l, \text{mid}]$ and $[\text{mid}+1, r]$ and sort them respectively.
    \item[]   3. Combine: Maintain two pointers $i$ and $j$ for the left and right subarrays.
                While merging, if $P[i] \leq P[j]$, then no new inversions are found, and we move the pointer $i$ backward.
                If $P[i] > P[j]$, then all remaining elements in the left subarray (from $i$ to $\text{mid}$) are greater than $P[j]$, which means there are $(\text{mid} - i + 1)$ inversions. Add this count to our total inversions and move the pointer $j$.
                Continue this process until all elements are merged, then return the total number of inversions found in this merge step.
\end{enumerate}

For every inversion $(p_i, p_j)$, given a division, they can be classified into three cases:

1. Both $p_i$ and $p_j$ are in the left subarray
2. Both $p_i$ and $p_j$ are in the right subarray
3. $p_i$ is in the left subarray and $p_j$ is in the right subarray

The first two cases are handled by the recursive calls, while the third case is handled during the merge step. Therefore, all inversions are counted exactly once.
\end{solution}

%---------------------------------------------------------
\part[3]

Analyze the time complexity of the complete algorithm.

\begin{solution}

$$
T(n) = 2T(n/2) + O(n)    
$$

$$
T(n) = O(n \log n)
$$

\end{solution}

\end{parts}
